% from ega4-21c-fr.tex
\oldpage[IV-4]{329}
(iv) Le résultat (21.14.3) a d'abord été démontré par C.~Seshadri [44] dans le cas
particulier où $S$ est un schéma algébrique normal sur un corps algébriquement clos
$k$ et $X=S\times_k T$, où $T$ est un $k$-préschéma algébrique lisse sur $k$. La
démonstration de Seshadri [44, p.~188--189] est de nature globale et utilise la
théorie des schémas de Picard. Elle donne en outre (\emph{loc. cit.}) des résultats
tels que le suivant (pour lequel on ne possède pas à l'heure actuelle de démonstration
par voie locale). Soient $S$, $T$ deux préschémas localement de type fini sur un
corps $k$, $X=S\times_k T$, $Z$ un cycle $1$-codimensionnel sur $X$ (considéré comme
$S$-préschéma); supposons vérifiées les conditions suivantes :
\begin{enumerate}[label=\arabic*\textsuperscript{o}]
  \item $S$ et $T$ sont géométriquement normaux sur $k$ (6.7.6);
  \item Pour tout point maximal $\eta$ de $S$, le cycle $1$-codimensionnel $Z_\eta$
  sur la fibre $X_\eta$, ayant même multiplicité que $Z$ en tout point de
  $X^{(1)}\cap X_\eta$, est localement principal (autrement dit, est l'image d'un
  diviseur de $X_\eta$, puisque $X_\eta$ est normal);
  \item Pour tout $s\in S$, $Z$ est principal aux points maximaux de la fibre $X_s$.
\end{enumerate}
Alors $Z$ est localement principal. En d'autres termes, $X$ étant normal (6.14.1),
pour tout $x\in X$ qui n'est pas maximal dans sa fibre $X_s$ et qui n'appartient à
aucune des « fibres génériques » $X_\eta$ (ce qui implique
$\dim(\mathcal O_{X,x})\geq2$ en vertu de (6.1.1)), l'anneau local
$\mathcal O_{X,x}$ est parafactoriel, en vertu de (21.12.15).
\footnote{En fait, dans l'article précité, Seshadri suppose que $k$ est algébriquement
clos, $T$ séparé et « semi-complet » (i.e.\ tel que
$\Gamma(T,\mathcal O_T)$ soit $k$-isomorphe à $k$) et remplace l'hypothèse
$3^{\mathrm o}$ par l'hypothèse plus forte que $\operatorname{Supp}(Z)$ ne contient
aucune des fibres $X_s$ pour $s\in S$. Mais comme l'énoncé est local sur $S$, on en
conclut aussitôt qu'il suffit de faire l'hypothèse $3^{\mathrm o}$, et cela prouve que
la conclusion (interprétée comme ci-dessus en termes de propriété de parafactorialité
des anneaux $\mathcal O_{X,x}$) est locale sur $S$ et sur $T$, ce qui permet
d'éliminer complètement l'hypothèse que $T$ est « semi-complet » et que $k$ est
algébriquement clos, puisque (quitte à passer d'abord à la clôture algébrique de $k$)
on peut supposer d'abord $T$ affine, ce qui permet de le plonger comme ouvert d'un
schéma projectif et normal sur $k$, auquel le résultat de Seshadri s'applique. Notons
aussi que, grâce à cette réduction, il suffit de faire la démonstration de Seshadri dans
le cas où $T$ est \emph{projectif} (et non seulement « semi-complet »), cas où la
théorie du schéma de Picard utilisée par Seshadri est contenue dans celle qui sera
développée au chap.~VI de notre Traité.}
\end{env}

\subsection{Diviseurs relatifs}
\EGAReaderTarget{ega4.c04.subsec.21-15}{subsection:IV.21.15}
\label{IV.21.15}

\begin{env}[21.15.1]
\label{IV.21.15.1}
Soient $S$ un préschéma, $f:X\to S$ un morphisme \emph{plat et localement de
présentation finie}. On a défini dans (20.6.1) le faisceau d'anneaux
$\mathscr M_{X/S}$ des \emph{germes de fonctions méromorphes sur $X$ relatives à
$S$}, sous-faisceau de $\mathscr M_X$; il est clair que l'injection canonique
$\mathcal O_X\to\mathscr M_X$ (20.1.4.1) applique $\mathcal O_X$ sur un
sous-faisceau de $\mathscr M_{X/S}$, avec lequel on l'identifie. Soit
$\mathscr M_{X/S}^*$ le faisceau (en groupes multiplicatifs) des germes de sections
inversibles de $\mathscr M_{X/S}$ qui est donc un sous-faisceau de $\mathscr M_X^*$
et contient $\mathcal O_X^*$ comme sous-faisceau.
\end{env}

\begin{definition}[21.15.2]
\EGAReaderTarget{ega4.c04.def.21-15-2}{IV.21.15.2}
On appelle \emph{faisceau des diviseurs sur $X$ relativement à $S$}, ou
\emph{faisceau des diviseurs sur $X$ transversaux à $f$}, et on note
$\mathscr{D}iv_{X/S}$ le faisceau quotient (de groupes commutatifs)
$\mathscr M_{X/S}^*/\mathcal O_X^*$; les sections de ce faisceau au-dessus de $X$ se
nomment \emph{diviseurs sur $X$ relatifs à $S$}, ou \emph{diviseurs sur $X$
transversaux à $f$}; ils forment un groupe commutatif noté $\operatorname{Div}(X/S)$.
\end{definition}

Il est clair que $\mathscr{D}iv_{X/S}$ s'identifie canoniquement à un sous-faisceau de
$\mathscr{D}iv_X$, et par suite $\operatorname{Div}(X/S)$ à un sous-groupe de
$\operatorname{Div}(X)$, que l'on note encore \emph{additivement}. Pour
